% AgonProject poster — AIG tikzposter style (A0 portrait)
% Screenshot-showcase layout. Style reused from the AIG sample poster;
% the previous sample content is preserved in sample_poster.tex.bak.

\documentclass[25pt,a0paper,portrait,margin=2mm, innermargin=16mm, blockverticalspace=4mm, colspace=15mm, subcolspace=10mm]{tikzposter}

\usepackage{aigpostersettings}

\usepackage{anyfontsize}
\usepackage{palatino}
\usepackage[utf8]{inputenc}
\usepackage[T1]{fontenc}
\usepackage{booktabs}
\usepackage[table]{xcolor}
\usepackage{graphicx}
\usepackage{amssymb,amsthm,amsmath,amsfonts}
\usepackage{qrcode}

\usepackage{xspace}

% Instance-specific project name macros
\newcommand{\agon}{\textsc{AgonProject}\xspace}
\newcommand{\tweety}{\textsc{TweetyProject}\xspace}

\usetitlestyle{aigpostertitle1}

\settitlefontsize{4cm}
\title{AgonProject}
\subtitle{An online platform for exploring approaches to formal argumentation}
%\author{Lars Bengel \quad Matti Berthold \quad Oleksandr Dzhychko \quad Matthias Thimm}
\institute{Artificial Intelligence Group, University of Hagen, Germany}

\usecolorstyle{aigpostercolor}
\usebackgroundstyle{aigposterback}
\useblockstyle{aigposterblockB}

\begin{document}

\maketitle

%======================================================================
% Row 1 — Motivation | QR + architecture
%======================================================================
\begin{columns}
\column{0.63}
\blueblock{Why AgonProject?}{
    Formal argumentation offers a rich \emph{zoo} of formalisms for reasoning with
    conflicting information --- but their diverse semantics, algorithms, and
    implementation details make them hard for students and newcomers to grasp.
    \tweety provides the most comprehensive collection of argumentation
    implementations available, yet exposes only \emph{text-based} interfaces for a
    handful of tasks.

    \vspace{0.4cm}
    \yellowcallout{%
    \textbf{Goal:} a unified, \textbf{visual}, install-free web interface to
    \emph{construct}, \emph{evaluate}, and \emph{compare} argumentation formalisms
    --- right in the browser.}
}
\column{0.37}
\yellowblock{Try it now --- no installation}{
    \centering
    {\setlength{\fboxsep}{10pt}\colorbox{white}{\qrcode[height=4.5cm]{http://agon.tweetyproject.org}}}\\[0.7ex]
    {\Large\bfseries\color{aigblue}\texttt{agon.tweetyproject.org}}
}
\blueblock{Built for exploration}{
    A client-side single-page app (Vue~3 + TypeScript) that delegates reasoning to a
    \tweety REST backend --- backend-agnostic, open source under GPL-3.0.
}
\end{columns}

%======================================================================
% Row 2 — Six formalisms | landing screenshot
%======================================================================
\begin{columns}
\column{0.58}
\blueblock{Six Argumentation Formalisms}{
    \begin{description}\setlength{\itemsep}{0.6ex}
        \item[AF] \emph{Abstract} argumentation --- arguments and directed attacks (Dung).
        \item[BAF] \emph{Bipolar} --- attacks \textbf{and} a support relation.
        \item[ADF] \emph{Abstract dialectical} --- propositional acceptance conditions.
        \item[iAF] \emph{Incomplete} --- uncertain arguments and attacks.
        \item[PAF] \emph{Probabilistic} --- probabilities on arguments and attacks.
        \item[SetAF] \emph{Collective attacks} --- sets of arguments attack a target.
    \end{description}
    \vspace{0.3ex}
    Every formalism shares the same editor, evaluation, export, and learning features.
}
\column{0.42}
\yellowblock{One Landing Page}{
    \shotpanel{figures/screenshot-landing.png}{Pick a formalism and start modelling
    straight away --- with curated examples and one-click random generation.}
}
\end{columns}

%======================================================================
% Row 3 — Five individual capability blocks
%======================================================================
\begin{columns}
\column{0.185}
\blueblock{Build}{Interactive graph editor with automatic layouts (Graphviz) and a
    physics simulation.}
\column{0.185}
\yellowblock{Evaluate}{Extension- and ranking-based semantics; credulous / sceptical
    acceptance, highlighted in the graph.}
\column{0.185}
\blueblock{Export \& Share}{\LaTeX{} (\texttt{argumentation}), SVG, TGF and ICCMA
    AF-format --- plus share-URLs.}
\column{0.185}
\yellowblock{Generate}{Random frameworks from eight configurable \texttt{networkX}
    generators.}
\column{0.185}
\blueblock{Learn}{10+ step-by-step tutorials and a glossary with literature references.}
\end{columns}

%======================================================================
% Row 4 — Three individual screenshot panels
%======================================================================
\begin{columns}
\column{0.315}
\yellowblock{Evaluate \& Rank (AF)}{
    \shotpanel{figures/screenshot-af.png}{Preferred extensions enumerated, with the
    \emph{categoriser} ranking annotated on each argument.}
}
\column{0.315}
\blueblock{Acceptance Conditions (ADF)}{
    \shotpanel{figures/screenshot-adf.png}{Propositional acceptance conditions,
    evaluated over three-valued interpretations; a stable model highlighted.}
}
\column{0.315}
\yellowblock{Collective Attacks \& Export (SetAF)}{
    \shotpanel{figures/screenshot-setaf.png}{Set-attacks drawn as hyperedges; the export
    window emits ready-to-use \LaTeX{} with an SVG preview.}
}
\end{columns}

%======================================================================
% Row 5 — Semantics + Outlook
%======================================================================
\begin{columns}
\column{0.58}
\block{Supported Semantics}{
    \small
    \renewcommand{\arraystretch}{1.25}
    \begin{tabular}{@{}p{0.30\linewidth}p{0.64\linewidth}@{}}
        \toprule
        \textbf{Group} & \textbf{Semantics}\\
        \midrule
        Classical & conflict-free, admissible, complete, grounded, preferred, stable\\
        Admissibility-based & ideal, semi-stable, eager, strong admissibility, initial sets, unchallenged\\
        Non-admissible & naive, CF2, SCF2, stage, stage2, (strongly) undisputed\\
        Weak-admissible & w-admissible, w-complete, w-grounded, w-preferred\\
        \midrule
        Rankings & burden, categoriser, counting, discussion, iterated graded defense, serialisability, social, strategy, tuples\\
        \bottomrule
    \end{tabular}

    \vspace{0.8ex}
    {\footnotesize Available across AF, BAF, iAF, PAF and SetAF; ADFs use the three-valued
    classical and naive semantics. Serialisation sequences are supported interactively.}
}
\column{0.42}
\useblockstyle{aigposterblockY}
\block{Outlook}{
    \vspace{0.2cm}
    Planned extensions:
    \begin{itemize}\setlength{\itemsep}{0.5ex}
        \item further formalisms --- Assumption-Based Argumentation, Weighted AFs, Bipolar
        Set-based AFs, claim-augmented AFs;
        \item benchmark generators from the ICCMA ecosystem;
        \item explanation-based reasoning: \emph{why} is an argument accepted?
    \end{itemize}
    \vspace{0.3cm}
    \yellowcallout{\centering
    \agon is backend-agnostic and open to new formalisms and solvers.}
}\useblockstyle{aigposterblockB}
\end{columns}

%======================================================================
% Footer
%======================================================================
\begin{columns}
\column{1.0}
\block{}{
    \centering
    {\large
    \textbf{Live:}~\texttt{agon.tweetyproject.org} \quad$\bullet$\quad
    \textbf{Source:}~\texttt{github.com/aig-hagen/AgonProject} \quad$\bullet$\quad
    GPL-3.0}\\[0.6ex]
    {\footnotesize Partially supported by the Deutsche Forschungsgemeinschaft (project
    ``Argumentative Reasoning in Nonsensical Situations'', ARNS, grant 550735820).}
}
\end{columns}

\end{document}
